% !TEX root = ../main.tex

% Local Variables:
% TeX-master: "../main"
% End:
% chktex-file 26

%%%%%%%%%%%%%%%%%%%%%%%%%%%%%%% Header %%%%%%%%%%%%%%%%%%%%%%%%%%%%%%%%%%%%%%%%%%%%
\begin{minipage}[l]{0.42\textwidth}
    \includegraphics[width=1\textwidth]{img/logo-UNAMBA.png}
\end{minipage}
\hfill
\begin{minipage}[c]{0.5\textwidth}
    \begin{flushright}
	\large{\textbf{Unidad \#2}}\\
	\large{Lectures on Física III}\\
	\large{24 de Julio del 2025. Haquira, Apurimac}\\
        % \large{\textbf{Student:} Huallpa Aimituma Josué David}
    \end{flushright}
\end{minipage}
\begin{tikzpicture}
    \draw[gray,thick] (-6.5,0)--(14,0);
\end{tikzpicture}


 %%%%%%%%%%%%%%%%%%%%%%%% INICIO DEL CONTENIDO EN DOS COLUMNAS %%%%%%%%%%%%%%%%%%%%%
  
 \begin{multicols}{2}
   \begin{center}
         \LARGE{\textbf{Capítulo IV: Circuitos de corriente continua}}\\	
         \vspace{0.2cm}
         % \Large {Lecturers Esteban Chalbaud \& Daniel Galviz} \\
         % \large{Teaching Assistant: Mauricio Gamonal \& Irvin Martínez}\\
         % \large{PhysicsLatam.com}\\
         % \vspace{0.2cm}
         \large{12 Agosto 2025, 5:59 pm (GMT-4)}\\
         % \vspace{0.2cm}
         \large{— Evaluación Parcial 2 —}
     \end{center}
    %%%%%%%%%%%%%%%%%%%%%%%%%%%%excercise%%%%%%%%%%%%%%%%%%%%%%%%%%%%%%%%%%%%%%%%
    \begin{excercise}[][][]{ex:k28}{(\textbf{4 pts})
        \textbf{(Preguntas teóricas)}
        \begin{itemize}
            \item[a)] Indique si la capacitancia de un condensador depende de la diferencia de potencial entre sus placas. Describa que pudiera suceder cuando esta diferencia de potencial crece inconmensurablemente.
            \item[b)] Cuando se conectan dos capacitores $C_1$ y $C_2$ en serie, la capacidad equivalente siempre es menor que $C_1$ o $C_2$. Por otra parte, si se conetan en paralelo la capacidad equivalente es mayor a $C_1$ o $C_2$. Explique estos hechos de forma cualitativa, a partir de un diagrama encuentre la relación entre la capacidad equivalente y las capacidades individuales de 3 condensadores conectados en serie y paralelo. 
            \item[c)] El máximo voltaje disponible para almacenar $1.2\times 10^{-5}$ C de carga de un capacitor  es de 1500 V. Determine la máxima capacidad necesaria.
            \item[d)] Describa dos funciones útiles que realiza un material aislante o dieléctrico, cuando se utiliza en un capacitor. 
        \end{itemize}
         }
    \end{excercise}
    %%%%%%%%%%%%%%%%%%%%%%%%%%%%excercise%%%%%%%%%%%%%%%%%%%%%%%%%%%%%%%%%%%%%%%%
    \begin{excercise}[][][]{ex:k29}{(\textbf{2 pts})
        Que área debe tener un condensador de placas paralelas separadas una distancia de 1cm para almacenar:
        \begin{itemize}
             \item[a)] 1F
             \item[b)] 1pF            
             \item[c)] $1\mu$ F
         \end{itemize}
         Para cada caso, parta de la definición de capacidad dado por:
         \begin{equation*}
             C=\frac{Q}{\Delta V}
         \end{equation*}
         }
    \end{excercise}
    \begin{excercise}[][][]{ex:k29}{(\textbf{4 pts})       Determine las fórmulas de la capacidad de condensadores de placas paralelas: 
        \begin{itemize}
             \item[a)] Planas separadas una distancia $d$
             \item[b)] Cilindricas de radio interior $a$, exterior $b$ y longitud $L$
             \item[c)] Esféricas de radio interior $a$, exterior $b$  
         \end{itemize}
         }
    \end{excercise}
    \begin{excercise}[][][$d=2.6\,\mathrm{m}$]{ex:11}{ (\textbf{3 pts})  
        Un capacitor tubular  consiste en dos cilindros conductores coaxiales muy largos. El radio externo del cilindro interior es de 9.50 cm y el radio interno del cilindro exterior es de 10.0 cm.
        \begin{itemize}
             \item[a)]  Determine la capacitancia de esta estructura, para una longitud de 30 cm de largo.. 
             \item[b)]  Obtenga la magnitud de la diferencia de potencial o voltaje que puede soportar esta estructura sin que sufra ruptura eléctrica, suponiendo que la disrupción eléctrica del aire situado en la región entre las placas, ocurre cuando se establece un campo eléctrico de 3.0 x 106 V/m en cualquier punto.
         \end{itemize}   
     }
    \end{excercise}

    %%%%%%%%%%%%%%%%%%%%%%%%%%%%excercise%%%%%%%%%%%%%%%%%%%%%%%%%%%%%%%%%%%%%%%%
    \begin{excercise}[][][]{ex:k30}{(\textbf{5.5 pts})
       Un condensador plano tiene placas cuadradas de 10cm de lado y una separación de d=4mm. Un bloque dieléctrico de k=2 tiene dimensiones de 10cm x 10cm x 4mm 
       \begin{itemize}
            \item[a)] Calcular la capacidad del condensador sin dielétrico
            \item[b)] ¿Cuál es la capacidad si el bloque dieléctrico llena el espacio entre las placas?
        \end{itemize} 
        }
    \end{excercise}
    %%%%%%%%%%%%%%%%%%%%%%%%%%%%excercise%%%%%%%%%%%%%%%%%%%%%%%%%%%%%%%%%%%%%%%%
    \begin{excercise}[][][]{ex:25}{(\textbf{2 pts})
       Determinar el valor del campo eléctrico en un cable de cobre de calibre 22 cuando este transporta una corriente de 5.2 A. (Diámetro nominal calibre 22: 0.6438 mm, $\sigma_{cu}=5.96\times 10^7\ Sm^{-1}$)
       } 
    \end{excercise}
    %%%%%%%%%%%%%%%%%%%%%%%%%%%excercise%%%%%%%%%%%%%%%%%%%%%%%%%%%%%%%%%%%%%%%%
     \begin{excercise}[][][]{ex:19}{(\textbf{3 pts})
        Para la Fig. Calcular La capacidad equivalente del sistema si cada capacitor es de $500\mu$ F   
         \begin{figure}[H]
             \centering
             \includegraphics[width=0.8\linewidth]{img/01_physics-i/03_statics/29_ex.png}
         \end{figure}
    } 
     \end{excercise}
    %%%%%%%%%%%%%%%%%%%%%%%%%%%%excercise%%%%%%%%%%%%%%%%%%%%%%%%%%%%%%%%%%%%%%%%
    \begin{excercise}[][][]{ex:19}{(\textbf{3 pts})
       Para la Fig. Calcular La resistencia equivalente si cada resistor es de 200 $\Omega$  
        \begin{figure}[H]
            \centering
            \includegraphics[width=\linewidth]{img/01_physics-i/03_statics/28.1.png}
        \end{figure}
    } 
    \end{excercise}
\end{multicols}
